\documentclass[11pt,a4paper]{article}

% --- Paquetes de Formato e Idioma ---
\usepackage[spanish,es-tabla]{babel}
\usepackage{fontspec}
\usepackage[margin=2.5cm]{geometry}
\usepackage{xcolor}
\usepackage{booktabs}
\usepackage{listings}
\usepackage{microtype}
\usepackage{metalogo}
\usepackage{amsmath}



% --- Configuración Tipográfica de la Documentación ---
\setmainfont{Arial}
\setsansfont{Noto Sans}

% --- Colores Institucionales Reutilizados de iesibai.sty ---
\definecolor{iesblue}{HTML}{0F2C59}
\definecolor{iesblueprimary}{HTML}{1A365D}
\definecolor{iesgray}{HTML}{4A5568}
\definecolor{iesaccent}{HTML}{DD6B20}

% --- Configuración de Hipervínculos ---
\usepackage[
    colorlinks=true,
    linkcolor=iesblueprimary,
    urlcolor=iesaccent,
    citecolor=iesgray,
    pdftitle={Documentación del paquete iesibai},
    pdfauthor={F. Javier Morales M. (TheBiotechScientist)}
]{hyperref}

% --- Estilo para Bloques de Código ---
\lstset{
    language=[LaTeX]TeX,
    basicstyle=\small\ttfamily,
    keywordstyle=\color{iesblueprimary}\bfseries,
    commentstyle=\color{iesgray}\itshape,
    stringstyle=\color{iesaccent},
    backgroundcolor=\color{gray!6},
    frame=single,
    framerule=0.4pt,
    rulecolor=\color{gray!30},
    breaklines=true,
    tabsize=2,
    showstringspaces=false
}

% --- Macros Auxiliares ---
\newcommand{\pkg}[1]{\textbf{#1}}
\newcommand{\cmd}[1]{\texttt{\textbackslash#1}}
\newcommand{\key}[1]{\texttt{#1}}

\title{%
    \Huge \textbf{\color{iesblueprimary}El paquete \pkg{iesibai}}\\[0.3em]
    \large Plantilla \LaTeX \ (No Oficial) para el \textbf{IESIBAI}\\[0.2em]
    \small Manual de Usuario y Especificación Técnica
}
\author{\textbf{F. Javier Morales M. \textit{(TheBiotechScientist)}}\\[10pt] {\bfseries \ttfamily \href{mailto:fmoralesm87@gmail.com}{fmoralesm87@gmail.com}}}

\date{\today}
\date{Versión 1.0.0 --- 27 de septiembre de 2026}

\begin{document}

\maketitle

\begin{abstract}
El paquete \pkg{iesibai} proporciona macros, estilos tipográficos, paletas de color y una interfaz clave-valor (\emph{key-value}) para estandarizar la presentación gráfica de trabajos académicos e institucionales en el \textbf{Instituto de Estudios Superiores de la Red Iberoamericana de Academias de Investigación (IESIBAI)}.
\end{abstract}

\tableofcontents
\vspace{0.8cm}
\hrule
\vspace{0.8cm}

\section{Introducción}
El objetivo principal de \pkg{iesibai} es automatizar la creación de portadas, la configuración de márgenes y el formato tipográfico de los documentos oficiales del instituto. Con este paquete, el autor puede concentrarse en la redacción del contenido mientras el sistema mantiene la coherencia gráfica requerida.

\section{Requisitos del Sistema y Compilación}
Este paquete hace uso intensivo de la gestión moderna de fuentes mediante \pkg{fontspec}. Por lo tanto, \textbf{es obligatorio compilar el documento con XeLaTeX}.

\subsection{Tipografías Requeridas}
El sistema operativo debe contar con las siguientes familias tipográficas instaladas:
\begin{itemize}
    \item \textbf{Arial}: Utilizada como fuente principal para el cuerpo del texto.
    \item \textbf{Helvetica}: Aplicada en títulos de portada y elementos destacados.
    \item \textbf{Noto Sans}: Utilizada como tipografía sin serifa auxiliar.
\end{itemize}

\subsection{Paquetes Dependientes}
El paquete \pkg{iesibai} carga automáticamente los siguientes paquetes:
\begin{itemize}
    \item \pkg{geometry}: Configurado para hoja tamaño carta (\texttt{letterpaper}) con márgenes de $2.5\ \text{cm}$.
    \item \pkg{babel}: Idioma español con soporte para tablas (\texttt{es-tabla}) y renombrado de índice a "`Contenido"'.
    \item \pkg{xcolor}, \pkg{graphicx}, \pkg{eso-pic}, \pkg{tikz}: Para gráficos, fondos de página completa y diseño visual.
    \item \pkg{titlesec}: Para la personalización de secciones y capítulos.
    \item \pkg{pgfkeys}, \pkg{hyperref}, \pkg{setspace}, \pkg{metalogo}.
\end{itemize}

\section{Comandos Principales}

\subsection{Configuración de Metadatos: \cmd{iesibaisetup}}
Asigna valores a las variables institucionales. Se recomienda colocar este comando en el preámbulo del documento.

\begin{lstlisting}
\iesibaisetup{
    institucion = {Instituto de Estudios Superiores de la Red Iberoamericana},
    campus      = {Campus Virtual Hidalgo},
    programa    = {Maestría en Educación},
    materia     = {Tecnología e Innovación I},
    subtitulo   = {Diseño de un paquete de entorno académico},
    trabajo     = {Proyecto Final},
    alumno      = {Nombre del Alumno},
    profesor    = {Dr. TheBiotechScientist},
    ubicacion   = {San Luis Potosí, México},
    fecha       = {Septiembre, 2026}
}
\end{lstlisting}

\subsection{Generación y Control de Portada}
Por defecto, \pkg{iesibai} \textbf{genera automáticamente la portada} al ejecutar \texttt{\textbackslash begin\{document\}} utilizando las variables declaradas.

\begin{itemize}
    \item \cmd{sinportadaiesibai}: Desactiva la generación automática de la portada. Debe colocarse en el preámbulo si se desea omitirla.
    \item \cmd{portadaiesibai}[opciones]: Permite compilar manualmente la portada en un punto específico del documento o pasarle opciones clave-valor directas.
\end{itemize}

\begin{lstlisting}
% Para desactivar la portada automática:
\sinportadaiesibai

% Para invocar la portada de forma manual con cambios puntuales:
\portadaiesibai[trabajo={Avance del Módulo II}]
\end{lstlisting}

\subsection{Lectura de Parámetros: \cmd{iesibaiget}}
Permite recuperar en cualquier parte del texto el valor asignado a una variable:

\begin{lstlisting}
El presente reporte corresponde a la materia \iesibaiget{materia}.
\end{lstlisting}

\section{Opciones de Claves (\emph{Keys})}

La Tabla~\ref{tab:claves} lista las claves disponibles en \cmd{iesibaisetup}, sus valores por defecto y su descripción.

\begin{table}[h]
\centering
\small
\begin{tabular}{lll}
\toprule
\textbf{Clave} & \textbf{Valor por Defecto} & \textbf{Descripción} \\
\midrule
\key{institucion} & \textit{Nombre del Instituto} & Nombre oficial de la institución. \\
\key{campus}      & \textit{Nombre/Tipo de Campus} & Sede o modalidad académica. \\
\key{programa}    & \textit{Maestría en Educación} & Programa o posgrado. \\
\key{materia}     & \textit{Titulo de la Materia}  & Nombre de la asignatura. \\
\key{subtitulo}   & \textit{Un subtitulo opcional...}& Subtítulo explicativo para la portada. \\
\key{trabajo}     & \textit{Trabajo del x módulo}  & Tipo de entrega o proyecto. \\
\key{alumno}      & \textit{Nombre del Alumno}     & Nombre del estudiante. \\
\key{profesor}    & \textit{Dr. Jorge Ruiz...}     & Nombre del docente/asesor. \\
\key{ubicacion}   & \textit{Estado, México}        & Ciudad y estado. \\
\key{fecha}       & \textit{Septiembre, 2026}      & Fecha del documento. \\
\key{engine}      & \textit{Elaborado en \XeLaTeX} & Nota sobre el motor de compilación. \\
\key{logo}        & \texttt{assets/logo\_iesibai.png} & Ruta del logotipo institucional. \\
\key{fondo}       & \texttt{assets/fondo\_...jpeg} & Ruta de la imagen de fondo. \\
\bottomrule
\end{tabular}
\caption{Claves configurables con el comando \cmd{iesibaisetup}.}
\label{tab:claves}
\end{table}

\section{Paleta de Colores Institucionales}
El paquete define cuatro colores globales en modelo HTML que pueden emplearse para personalizar tablas o cajas de texto en el documento:

\begin{itemize}
    \item \textbf{iesblueprimary} (\texttt{\#1A365D}): Azul marino profundo principal (utilizado en títulos y acentos institucionales).
    \item \textbf{iesblue} (\texttt{\#0F2C59}): Azul oscuro complementario.
    \item \textbf{iesaccent} (\texttt{\#DD6B20}): Naranja/coral sutil para resaltar subtítulos y enlaces.
    \item \textbf{iesgray} (\texttt{\#4A5568}): Gris oscuro para textos secundarios y detalles de portada.
\end{itemize}

\section{Documento de Ejemplo Mínimo}

A continuación se presenta un código funcional completo:

\begin{lstlisting}
\documentclass[12pt]{article}
\usepackage{iesibai}

\iesibaisetup{
    materia   = {Tecnología e Innovación I},
    subtitulo = {Implementación de Entornos LaTeX},
    alumno    = {TheBiotechScientist},
    profesor  = {Dr. Jorge Ruiz Mondragón},
    fecha     = {Septiembre, 2026}
}

\begin{document}
% La portada se inserta automáticamente al iniciar el documento.

\section{Introducción}
Este es el inicio del cuerpo del documento redactado con la tipografía y márgenes estandarizados de IESIBAI.

\end{document}
\end{lstlisting}

\section{Licencia}
Este paquete se distribuye bajo la licencia \textbf{\LaTeX\ Project Public License (LPPL)} versión 1.3c o posterior. Para más detalles, consulte: \url{https://www.latex-project.org/lppl.txt}.

\end{document}